\section{Path stability and the size of a marked subgraph}
\label{sec:research-path}

These results use the forest proof in the November 10 manuscript.
The parameter extension below is a mathematical argument, not a Lean formalization.
It first improves the general path bound to order $b/\log b$.
Theorem~\ref{thm:general-linear-path-stability} later obtains the full linear bound from this enlarged forest range.
Theorem~\ref{thm:general-additive-path-stability} then strengthens it to an additive loss depending only on the density gap.
For two-terminal series-parallel networks, an elementary argument gives an explicit logarithmic loss.

All logarithms in this section have base two, except for $\ln$ and $\exp$.
Write $P_{uv}(N)$ for the number of simple paths between distinct vertices $u,v$.
Write $b=\beta(N)$.

\subsection{From paths to forest outside parts}

\begin{lemma}\label{lem:research-path-injection}
Let $N$ be a connected simple graph, and let $A$ be any $u$--$v$ path of length $\ell$.
Add a new vertex $w$ and edges $uw,wv$ to obtain $H$.
Let $W=A\cup\{uw,wv\}$.
Then
\[
 \beta(H)=b+1,\qquad \beta(W)=1,\qquad |E(W)|=\ell+2,
 \qquad P_{uv}(N)\le p(H,W)2^b.
\]
\end{lemma}

\begin{proof}
Each terminal path $P$ gives the cycle $P\cup\{uw,wv\}$ of $H$.
Its outside part is $P\setminus E(A)$, which is a forest.
Suppose that $P\setminus E(A)=Q\setminus E(A)$.
Then $P\mathbin\triangle Q$ is an even edge set supported on the forest $A$.
It must be empty, so $P=Q$.
There are $2^{\beta(H)-\beta(W)}=2^b$ outside parts.
Their forest proportion is $p(H,W)$.
\end{proof}

The formalized forest theorem gives $\ell>\sqrt{b+1}-2$ for fixed positive path-density gap and sufficiently large $b$.
The manuscript's size extension gives $\ell>(b+1)/\log^2(b+1)-2$.
The next argument removes one logarithmic factor for a fixed gap.

\subsection{A fixed-gap forest estimate at size \texorpdfstring{$r/\log r$}{r/log r}}

\begin{theorem}\label{thm:research-marked-size}
Fix $0<\delta<1/2$, and set
\[
 A_\delta=6\exp(16/\delta),\qquad
 a_\delta=\frac{1}{12A_\delta}.
\]
For sufficiently large $r$, the following statement holds.
Let $G$ be a connected simple graph of rank $r$.
Let $W$ be a nonempty union of cycles satisfying
\[
 |E(W)|\le a_\delta\frac{r}{\log r},\qquad
 \operatorname{comp}(W)\le\frac{\log r}{100}.
\]
Then $p(G,W)\le1/2+\delta$.
\end{theorem}

\begin{proof}
We use the manuscript's reduction, cut lemmas, spectral estimate, and second-moment argument.
The details below check each estimate affected by the new parameters.
Constants in $O_\delta(\cdot)$ may depend on $\delta$.

Put $x=\log r$, $y=\log x$, $D=\lfloor x/10\rfloor$, and $s=\lfloor(D-1)/2\rfloor$.
Set $L=\lfloor A_\delta x\rfloor$.
The reduction gives a subcubic multigraph $\Gamma$ and a marked vertex set $P$.
Every vertex outside $P$ has degree three, $|P|\le2|E(W)|$, and $\Gamma[P]$ has at most $x/100$ components.
Also $\beta(\Gamma)=r$ and $2r-2\le |V(\Gamma)|\le2r+2|E(W)|$.
Let $\mathcal F$ mean that the random even subgraph is a forest on $\Gamma-P$.
It suffices to rule out $\mathbb P(\mathcal F)>1/2+\delta$.

Assume this strict inequality.
Choose a maximal vertex-disjoint family $\mathcal Q$ of cycles of length at most $D$ in $\Gamma-P$.
The small-cut lemma gives
\[
 |\mathcal Q|<\frac{(D+1)2^{D-1}}{\delta}
       =O_\delta(xr^{1/10}).
\]
For their vertex union $S$, this implies
\[
 |S|=O_\delta(x^2r^{1/10})=o(r/x).
\]
Let $J=\Gamma-(P\cup S)$, and put $v=|V(J)|$.
This graph is simple and has girth greater than $D$.
The vertex and deficit estimates become
\[
 v=2r+O_\delta(r/x),\qquad
 d_J:=\sum_{z\in V(J)}(3-\deg_Jz)
 \le 3(|P|+|S|)
 \le\frac{6a_\delta r}{x}+o(r/x).
\]
Consequently,
\[
 \frac{Ld_J}{2v}\le\frac32 A_\delta a_\delta+o(1)
                 =\frac18+o(1).
\]

Let $B$ be the non-backtracking matrix of $J$.
For even $j$ with $6x\le j\le L$, the spectral estimate gives
\[
 2^{-j}\operatorname{tr}B^j
 \ge\left(1-\frac{d_J}{2v}\right)^j-3v2^{-j/2}
 \ge1-\frac{Ld_J}{2v}-3vr^{-3}\ge\frac12
\]
for sufficiently large $r$.
For the cycles $C$ of $J$ of length at most $L$, put
$\mu=\sum_C2^{-|C|}$.
The repeated-vertex walk estimate from the manuscript has error at most
\[
 \frac32 L^2 2^{-s}\mu=o(\mu).
\]
Therefore
\[
 (1+o(1))\mu
 \ge\frac12\sum_{\substack{6x\le j\le L\\j\ \mathrm{even}}}\frac1{2j}
 =\frac18\ln(A_\delta/6)+o(1).
\]
In particular, $\mu\ge\frac1{16}\ln(A_\delta/6)=1/\delta$ for sufficiently large $r$.

It remains to check the pair correlations.
Use $t_0=\lfloor y/2\rfloor$ in the manuscript's link argument.
For disjoint cycles $C,C'$ of length at most $L$, it gives
\[
 \ell(C,C')\le \operatorname{comp}(\Gamma[P])
 +\frac{2L}{t_0}
 +\frac{(t_0+1)2^{t_0-1}}{\delta}
 +\left(\frac{8L}{D}\right)^2
 \le\frac{x}{100}+o(x).
\]
Here the three added terms are respectively $O_\delta(x/y)$, $O_\delta(y\sqrt{x})$, and $O_\delta(1)$.
The exceptional-partner estimate thus remains
\[
 2^{x/100+o(x)}L^3 2^{-s}
 \le2^{x/100-x/20+o(x)}=o(1).
\]
The same weighted variable $Z$ therefore satisfies
\[
 \mathbb EZ=\mu,\qquad
 \mathbb EZ^2\le2\mu^2+(1+o(1))\mu.
\]
Since $\mathcal F$ implies $Z=0$, the second-moment lemma yields
\[
 \mathbb P(\mathcal F)\le\frac12+\frac{1+o(1)}{4\mu}
 \le\frac12+\frac\delta2
\]
for sufficiently large $r$.
This contradicts the assumption.
\end{proof}

\begin{corollary}\label{cor:research-general-path}
Fix $0<\varepsilon\le1/2$.
For sufficiently large $b$, if
\[
 P_{uv}(N)\ge(1/2+\varepsilon)2^b,
\]
then every terminal path has length
\[
 \ell>a_{\varepsilon/2}\frac{b+1}{\log(b+1)}-2.
\]
In particular, $\ell=\Omega_\varepsilon(b/\log b)$.
\end{corollary}

\begin{proof}
Apply Lemma~\ref{lem:research-path-injection}.
If the displayed length inequality failed, Theorem~\ref{thm:research-marked-size} would give
$P_{uv}(N)\le(1/2+\varepsilon/2)2^b$, a contradiction.
The one-component condition holds for sufficiently large $b$.
\end{proof}

The threshold is not explicit because the supporting asymptotic estimates are not explicit.
The additive $-2$ comes from the added two-edge path.
Neither issue affects the asymptotic order.

\begin{corollary}\label{cor:research-small-marked}
Suppose $r\to\infty$ along connected simple graphs $G$ and nonempty unions of cycles $W$.
If $|E(W)|=o(r/\log r)$ and $\operatorname{comp}(W)\le(\log r)/100$, then
$p(G,W)\le1/2+o(1)$.
\end{corollary}

\begin{proof}
For each fixed $0<\delta<1/2$, the size hypothesis eventually meets Theorem~\ref{thm:research-marked-size}.
Thus $\limsup p(G,W)\le1/2$.
\end{proof}

\subsection{An additive theorem for two-terminal series-parallel networks}

A two-terminal series-parallel network starts with one edge and its two endpoints.
Series composition identifies one terminal from each network.
Parallel composition identifies the corresponding terminal pairs.
Apart from those identifications, the two networks are disjoint.
Parallel edges are permitted and are counted separately.
This definition imposes a condition on the terminals, not only on the underlying graph.

\begin{theorem}\label{thm:research-ttsp}
Let $(N,u,v)$ be a two-terminal series-parallel network of rank $b$ and terminal distance $\ell$.
Then
\[
 P_{uv}(N)\le 2^{b-1}+2^{\ell-1}.
\]
Consequently, for $0<\varepsilon\le1/2$,
\[
 P_{uv}(N)\ge(1/2+\varepsilon)2^b
 \quad\Longrightarrow\quad
 \ell\ge b-\log\frac1{2\varepsilon}
 \quad\Longrightarrow\quad
 \ell\ge\frac{b}{\log(1/\varepsilon)}.
\]
\end{theorem}

\begin{proof}
Every terminal path is a join, so $P_{uv}(N)\le2^b$ for every connected network.
Define $h(t)=(1+2^{-t})/2$ for $t\ge0$.
We prove the stronger normalized statement
\[
 \frac{P_{uv}(N)}{2^b}\le h((b-\ell)_+),
 \qquad (z)_+=\max\{z,0\}.
\]
It holds for a single edge.

For series composition, write $b_i,\ell_i,P_i$ for the two input networks.
Then $b=b_1+b_2$, $\ell=\ell_1+\ell_2$, and $P=P_1P_2$.
For $s,t\ge0$,
\[
 h(s+t)-h(s)h(t)=\frac{(1-2^{-s})(1-2^{-t})}{4}\ge0.
\]
Also $(b_1-\ell_1)_++(b_2-\ell_2)_+\ge(b-\ell)_+$, and $h$ is decreasing.
The normalized statement follows by multiplication.

For parallel composition, $b=b_1+b_2+1$, $\ell=\min\{\ell_1,\ell_2\}\ge1$, and $P=P_1+P_2$.
Since $b_1,b_2\ge0$,
\[
 P\le2^{b_1}+2^{b_2}
 \le2^{b_1+b_2}+1
 =2^{b-1}+1
 \le2^{b-1}+2^{\ell-1}.
\]
Together with $P\le2^b$, this is the normalized statement.
Induction proves the theorem's first inequality.

The density hypothesis now gives $\varepsilon2^b\le2^{\ell-1}$.
Thus $b\le\ell+\log(1/(2\varepsilon))$.
Since $\ell\ge1$, the right side is at most $\ell\log(1/\varepsilon)$.
\end{proof}

The coefficient is sharp when $\varepsilon=2^{-b}$.
Take a series of $b-1$ triangles between $u,v$, and add the edge $uv$.
For $b\ge3$, the result is simple, has rank $b$, distance one, and $2^{b-1}+1$ terminal paths.
For $b=2$, a direct edge in parallel with two internally disjoint two-edge paths has the same count.
For $b=1$, use a triangle with adjacent terminals.

The theorem fails for arbitrary choices of terminals in a series-parallel graph.
On vertices $0,1,2,3,4$, take
\[
 E=\{01,02,03,04,12,14,23\},\qquad u=3,\quad v=4.
\]
This graph has rank three and terminal distance two.
Its seven terminal paths are
\[
 3014,\quad30214,\quad304,\quad32014,\quad3204,\quad32104,\quad3214.
\]
Thus $P=7>2^{3-1}+2^{2-1}=6$.
The graph is a triangle with two vertices successively attached to existing edges, so it has no $K_4$ minor.
It does not satisfy the two-terminal composition hypothesis for terminals $3,4$.

\subsection{An elementary structural reduction}

\begin{proposition}\label{prop:research-offroute}
Let $N$ be connected, with distinct terminals $u,v$.
Let $N_0$ be the union of the blocks on the $u$--$v$ route in the block-cut tree.
Count bridges as blocks.
Put $s=\beta(N)-\beta(N_0)$.
Then
\[
 P_{uv}(N)=P_{uv}(N_0)\le2^{\beta(N_0)}=2^{b-s}.
\]
In particular, $P_{uv}(N)>2^{b-1}$ implies $s=0$.
\end{proposition}

\begin{proof}
A simple terminal path cannot enter a block outside the terminal route.
It would have to revisit the articulation vertex where it left that route.
Thus all terminal paths lie in $N_0$.
Apply the join bound to $N_0$.
Ranks add across blocks, so $s$ is a nonnegative integer.
\end{proof}

This removes cyclic blocks outside the terminal route from the unresolved problem.
It does not control the internal structure of a single block.
The fixed-gap argument still loses a factor $\log b$ there.
Its spectral estimate counts walks of length at least order $\log b$.
A marked set of order $b$ creates a constant degree deficit, whose effect accumulates along those walks.
The present retuning does not prove general linear stability.

\paragraph{Finite checks.}
The accompanying script checks 1,728 distinct composition states for networks with at most 12 edges.
It verifies the inequality with exact integers across 27,264 series or parallel transitions.
It also enumerates the seven counterexample paths and the sharp simple examples for ranks three through nine.
These checks do not verify the asymptotic estimates or replace their proofs.
